\appendix

\section{Supplementary Figures}
\label{sec:appendix}

\begin{figure}[htbp]
\centering
\includegraphics[width=\textwidth]{v2/fig1_radar_behavioral_profile.pdf}
\caption{Multi-dimensional behavioral profiles by operating mode (TUB, STD, ECO). Values are normalized to [0,1] range for each outcome. TUB mode dominates on harsh events (aggressive riding) while ECO dominates on cruise fraction (stable riding). Speed CV is nearly identical across modes.}
\label{fig:radar}
\end{figure}

\begin{figure}[htbp]
\centering
\includegraphics[width=\textwidth]{v2/fig2_forest_gee_coefficients.pdf}
\caption{GEE coefficient forest plot for all primary outcomes (200,000-trip subsample). TUB is the reference mode. Points show estimated coefficients; error bars show 95\% confidence intervals. Safety outcomes (harsh events) show the strongest mode differentiation.}
\label{fig:forest}
\end{figure}

\begin{figure}[htbp]
\centering
\includegraphics[width=0.9\textwidth]{v2/fig3_shap_safety_harsh_event.pdf}
\caption{SHAP summary plot for the harsh event classification model (LightGBM, AUC = 0.734). TUB mode indicator dominates with mean $|\text{SHAP}| = 0.49$, confirming that mode selection is the primary determinant of harsh event occurrence.}
\label{fig:shap}
\end{figure}

\begin{figure}[htbp]
\centering
\includegraphics[width=\textwidth]{v2/fig4_did_multi_outcome.pdf}
\caption{Multi-outcome DiD coefficients with Bonferroni correction (KEY FIGURE). Red bars indicate Bonferroni-significant effects; orange bars indicate nominal significance ($p < 0.05$); gray bars indicate non-significant effects. Only harsh acceleration and deceleration survive Bonferroni correction.}
\label{fig:did_coef}
\end{figure}

\begin{figure}[htbp]
\centering
\includegraphics[width=\textwidth]{v2/fig5_demand_response.pdf}
\caption{Demand response to TUB ban by city treatment intensity. Trip counts and active users indexed to February 2023 baseline, stratified by TUB share quartiles. The December decline is seasonal (all quartiles drop similarly), not treatment-driven.}
\label{fig:demand}
\end{figure}

\begin{figure}[htbp]
\centering
\includegraphics[width=\textwidth]{v2/fig6_compensation_cohens_d.pdf}
\caption{Compensation test: Cohen's $d$ effect sizes for mode-switcher within-user DiD across all behavioral margins. All effects are negligible ($|d| < 0.2$), indicating zero behavioral compensation on any unconstrained margin.}
\label{fig:compensation}
\end{figure}

\begin{figure}[htbp]
\centering
\includegraphics[width=\textwidth]{v2/fig7_experience_trajectories.pdf}
\caption{Experience-mediated behavioral escalation. Left panel: TUB adoption rate by trip number. Right panels: outcome trajectories for all modes (dark line) vs.\ STD/ECO only (blue dashed). The gap between lines quantifies TUB-mediated effects: 57--70\% for safety outcomes, 18--27\% for dynamics.}
\label{fig:experience}
\end{figure}

\begin{figure}[htbp]
\centering
\includegraphics[width=\textwidth]{v2/fig8_survival_multi_endpoint.pdf}
\caption{Multi-endpoint Kaplan-Meier survival curves stratified by TUB usage. Left: first harsh event. Center: first high-CV trip (note: TUB is protective due to mechanical speed capping reducing CV). Right: first speeding event (non-TUB users median never reached).}
\label{fig:survival}
\end{figure}

\section{Event Study Details}
\label{sec:appendix_event}

\begin{figure}[htbp]
\centering
\includegraphics[width=\textwidth]{v2/fig_supp_event_study.pdf}
\caption{Event study: monthly treatment-intensity interaction coefficients for each outcome (reference: November 2023). The vertical red line indicates the TUB ban (December). Pre-period coefficients are generally close to zero; the December coefficient is the treatment effect.}
\label{fig:event_study}
\end{figure}

\section{Distributional Treatment Effects}
\label{sec:appendix_dist_twfe}

To examine whether the TUB ban's causal effects are concentrated at particular points of the outcome distribution, we extend the TWFE DiD (Equation~\ref{eq:twfe}) to city-month quantile summaries. For each city $c$ and month $t$, we compute the 25th, 50th, 75th, and 90th percentiles of each outcome across all trips, then estimate separate TWFE models at each quantile:
\begin{equation}
    Q_{ct}^{(j,\tau)} = \alpha_c^{(\tau)} + \gamma_t^{(\tau)} + \beta^{(j,\tau)} \cdot (\text{Post}_t \times \text{TUB}^{\text{Nov}}_c) + \epsilon_{ct}^{(\tau)}
    \label{eq:dist_twfe}
\end{equation}
where $Q_{ct}^{(j,\tau)}$ is the $\tau$-th percentile of outcome $j$ in city $c$ at month $t$, with $\tau \in \{0.25, 0.50, 0.75, 0.90\}$. Standard errors are clustered by city.

This approach estimates treatment effects at different points of the \textit{city-level} outcome distribution. It is \textit{not} a formal unconditional quantile treatment effect (QTE) estimator in the sense of \citet{firpo2009unconditional}, which requires individual-level data and recentered influence function (RIF) regressions. The city-level aggregation provides a descriptive complement to the mean-level DiD by revealing distributional heterogeneity in treatment effects.

Figure~\ref{fig:dist_twfe} presents the results. Mean speed exhibits strong distributional heterogeneity: the TUB ban reduces the 90th percentile by $-3.08$~km/h ($p < 10^{-16}$) but affects the 25th percentile by only $-0.75$~km/h (not significant, $p = 0.31$). The effect magnitude increases monotonically across quantiles ($\hat{\beta}_{p25} = -0.75^{\text{ns}}$, $\hat{\beta}_{p50} = -1.22^{*}$, $\hat{\beta}_{p75} = -1.93^{***}$, $\hat{\beta}_{p90} = -3.08^{***}$), though the lowest quantile is not individually significant. This gradient indicates that speed governors disproportionately compress the dangerous right tail of the speed distribution.

Cruise fraction reveals a hidden tail effect: the mean treatment effect is indistinguishable from zero ($\hat{\beta} = 0.04$, $p = 0.17$), but the 90th percentile shows a substantial positive shift ($\hat{\beta}_{p90} = +0.36$, $p < 10^{-7}$). The TUB ban forces high-speed trips into governed modes, dramatically increasing cruising behavior at the right tail.

Speed CV and zero-speed fraction are null at all quantiles, confirming the mean-level findings. Harsh event counts are ill-suited for distributional analysis because the city-month distributions are highly zero-inflated: the 25th and 50th percentiles are identically zero across virtually all cities and months, leaving no variation for the TWFE estimator at these quantiles. The mean-level effects nevertheless remain significant ($p < 0.001$ for both harsh acceleration and deceleration), as these are driven by changes in the upper tail of the count distribution.

\begin{figure}[htbp]
\centering
\includegraphics[width=\textwidth]{v2/fig_qte_exploration.pdf}
\caption{Distributional TWFE DiD: treatment effect coefficients by quantile for six behavioral outcomes. Each panel shows the TWFE DiD coefficient ($\hat{\beta}^{(j,\tau)}$) at the mean and four quantiles (p25, p50, p75, p90) of the city-month distribution. Error bars show 95\% confidence intervals (cluster-robust SEs). Mean speed shows effects increasing monotonically in magnitude across upper quantiles; cruise fraction reveals a large positive effect at the 90th percentile despite a null mean effect.}
\label{fig:dist_twfe}
\end{figure}

\section{Outcome Distributions}
\label{sec:appendix_dist}

\begin{figure}[htbp]
\centering
\includegraphics[width=\textwidth]{v2/fig_supp_distributions.pdf}
\caption{Violin plots of outcome distributions by operating mode (50,000 trips per mode). Black line: mean; white line: median. Harsh event distributions are highly right-skewed; cruise fraction and zero-speed fraction are approximately beta-distributed.}
\label{fig:distributions}
\end{figure}
